\section{From theory to circuits: classical reduction and the residual quantum core}\label{sec:reduction}
Feasibility by construction is only useful if the circuit can be run. We test this on a multi-period long/short portfolio with sector caps, long/short exclusivity, a per-period capital and budget, and a global trading budget: $A=12$ assets in three sectors and $T=8$ periods give $n=192$ variables and $137$ constraints of four kinds, with a quadratic risk objective (\cref{sec:large}).

\paragraph{The coherent state does not fit.} The automaton over all constraints has width $11$ at $n=24$ and $103$ at $n=192$, because the global budget couples all periods, and the compiled $A_q$ grows from $12.9$ thousand CZ gates and $88$ qubits to $5.1$ million and $1156$ qubits (default synthesis, before routing). The feasible share of the uniform state is $2^{-131}$, which is what the construction repairs, but the circuit is far beyond any device.

\paragraph{Reduction derived from the automaton.} Three classical steps act before any circuit is built. (i)~\emph{Variable order}: period-major gives widths $8$, $23$, $11$ against $36$, $172$, $36$ for asset-major on three small instances. (ii)~\emph{Conditioning on an interface}: drawing the per-period and per-sector counts of longs and shorts classically makes the periods independent blocks and the cost linear in $T$, from $5.1$ million to $29$ thousand CZ gates at $n=192$, and to about $2.9$ thousand with a second level of conditioning and one-hot registers. (iii)~\emph{Symmetry-aware lowering}: in each block the (long, short) pairs of the assets are interchangeable (an automatic check on $F$), so the block is prepared with Dicke states and a chain of controlled swaps, exactly uniform on the feasible set, with no register and no uncomputation; per block it is $3.1$--$4.4\times$ smaller than the generic circuit and $5$--$8\times$ when the generic registers must be uncomputed, as QAOA with pair-exchange mixers requires (for the preparation of a whole core the ratio is $14$--$28\times$, \cref{sec:sota}). This lowering is not specific to our construction (Dicke states are known); the role of the automaton there is to detect the structure and guarantee the support.

\paragraph{What conditioning costs.} The interface is drawn classically, so the superposition over interface values becomes a mixture. That is valid for sampling and for QAOA with mixers that conserve the interface (case (b) of \cref{thm:qaoa}), and invalid for amplification of the whole state and for the global Grover mixer, which need coherence across interface values.

\paragraph{Exact pruning, and how much is left.} Given the interface, variables that take the same value in all feasible strings of their block are known: they leave the objective, blocks with a single feasible string need no mixer, blocks without longs or shorts use a two-qubit swap rotation instead of a four-qubit pair exchange, and a constraint gauge rewrites couplings with the constant counts. The state on the feasible set is unchanged (checked against dense simulation to $10^{-14}$) and the compiled layer shrinks $4$--$8\times$ (at $n=24$, $1131\to259$--$262$ CZ, median ESP $0.01\to0.36$--$0.38$; \cref{tab:core_prune}). The reduction is large because the interface fixes most variables: on average only $2.5$, $6.0$, $6.9$ and $11.1$ of $8$, $12$, $16$ and $24$ variables stay free. This must qualify every size quoted for executability: what is left to the quantum computer is a \emph{core}, and its size, not that of the original problem, is what a device must run.

\begin{table}[t]
\centering\footnotesize\setlength{\tabcolsep}{4pt}
\caption{\orig{exact classical counting and sampling; no circuits} Residual quantum core after the exact classical reduction, over $200$ sampled interface tuples per row. Families differ in how loose the local constraints are (tight: at most one long position per sector and $3$ positions per period; medium: $2$, $6$; loose: $4$, $12$). Free: qubits that remain (upper bound); $\log_2|F|$: feasible strings of the core; ZZ: couplings among free variables for a dense objective.}\label{tab:core_map}
\begin{tabular}{lrrrr}
\toprule
family & $n$ & free & $\log_2|F|$ & ZZ\\
\midrule
tight & 96 & 42 & 18 & 881\\
tight & 192 & 84 & 25 & 3528\\
tight & 384 & 168 & 50 & 14112\\
tight & 768 & 336 & 64 & 56448\\
medium & 96 & 58 & 28 & 1691\\
medium & 192 & 113 & 44 & 6432\\
medium & 384 & 221 & 89 & 24605\\
medium & 768 & 437 & 120 & 95816\\
loose & 96 & 84 & 41 & 3535\\
loose & 192 & 137 & 71 & 9497\\
loose & 384 & 254 & 138 & 32406\\
loose & 768 & 478 & 203 & 114381\\
\bottomrule
\end{tabular}
\end{table}


\paragraph{How large can the core be?} The average above is for $n\le24$. \Cref{tab:core_map} counts the core for families of different looseness and up to $768$ variables: the share of free variables is $44$--$46\,\%$ for tight constraints and $62$--$88\,\%$ for loose ones, so the core grows with the problem, to $42$--$84$ free qubits for tight constraints at $n=96$--$192$ and to hundreds for loose ones. With tight constraints the core is still enumerable up to about $n=192$ ($2^{25}$ feasible strings) and needs no quantum treatment there; with loose constraints it has $2^{41}$ feasible strings at $n=96$ and $2^{138}$ at $n=384$, so enumeration is impossible and a quantum treatment of the core could matter. Such cores are far beyond what a device is predicted to run (\cref{sec:sota}). The reduction is preprocessing that benefits any quantum method equally; the size of the core is a property of how loose the constraints are.
